\documentclass[11pt]{article}
\usepackage{amsmath,amssymb,amsthm,hyperref}
\newtheorem{theorem}{Theorem}
\newtheorem{lemma}{Lemma}
\newcommand{\Q}{\mathbb Q}
\newcommand{\Z}{\mathbb Z}
\newcommand{\Fp}{\mathbb F_p}
\title{A computer-assisted exclusion at the quadratic regular rate}
\author{Research notes on Erd\H{o}s 243; no priority claim}
\date{6 November 2026}
\begin{document}
\maketitle

\noindent\textbf{Status.} This does not solve Erd\H{o}s 243. The final
irrationality theorem assumes the much stronger rate
$a_n^2/a_{n+1}=1+2/n+o(n^{-2})$. The unrestricted limit $1$ does not
imply this rate. No counterexample to the full problem is constructed.
The proof below uses two complete integral-point computations in Magma;
it does not use GRH or BSD as assumptions. Their input, proved full-group
flags, and output are retained separately for audit.

The exact rational-tail framework is from
\href{https://math.colgate.edu/~integers/aa28/aa28.pdf}{Koizumi (2026)}.
Gcd stabilization, adjacent coprimality, and Chebotarev square
specialisation are used in
\href{https://github.com/wcook04/plectis-erdos/blob/main/paper/243/erdos243-reciprocal-tail-reasoning-surface.tex}{Cook's companion paper}.
No novelty or priority claim is made here.

\paragraph{Integer orbits.}
Throughout an orbit means positive integers satisfying
\begin{equation}\label{eq:orbit}
 C_{n+1}=a_n C_n-D_n>0,\qquad D_{n+1}=a_nD_n.
\end{equation}
The same equations hold for $U_n=C_n/g$, $V_n=D_n/g$ if a fixed
integer $g$ divides both states. Two useful exact identities are
\begin{align}
 U_{n+2}&=(a_n+a_{n+1})U_{n+1}-a_n^2U_n,\label{eq:three}\\
 U_n V_{n+1}&=V_n(V_n+U_{n+1}).\label{eq:quadratic}
\end{align}
In particular, if $U_n\equiv0\pmod p$, then
\begin{equation}\label{eq:neighbours}
 -U_{n-1}U_{n+1}\equiv(a_{n-1}U_{n-1})^2\pmod p.
\end{equation}

\begin{lemma}[Two quadratic tests via Chebotarev]\label{lem:cheb}
Let $K=\Q(\alpha)$ be a number field. If $h\in K^\times$ is a
nonsquare, infinitely many rational primes have a reduction of $\alpha$
in $\Fp$ at which the reduction of $h$ is a nonsquare. If $h_1,h_2$
are both nonsquares in $K^\times$, infinitely many rational primes
have such a reduction at which both are nonsquares. In either statement
one may exclude any fixed finite set of bad primes.
\end{lemma}
\begin{proof}
In the first case the nontrivial automorphism of $K(\sqrt h)/K$
fixes $\alpha$ and negates $\sqrt h$. In the second case there is an
automorphism of $K(\sqrt{h_1},\sqrt{h_2})/K$ negating both square
roots: the two quadratic extensions are either equal or independent.
Extend that automorphism to a finite Galois extension $L/\Q$.
Chebotarev supplies infinitely many rational primes with its conjugacy
class as Frobenius; choose an above-prime having the specified
automorphism itself as Frobenius. Apart from the finitely many
denominator, ramification, characteristic-two, and zero-reduction
exceptions, $\overline\alpha$ is Frobenius-fixed and lies in $\Fp$,
whereas neither square root of each $\overline h_i$ is fixed.
Thus the indicated reductions are nonsquares. See
\href{https://www.jmilne.org/math/CourseNotes/CFT.pdf}{Milne,
Class Field Theory, Chapter V, Theorem 3.23} for the Chebotarev input.
\end{proof}

\begin{lemma}[Complete integral-point computations]\label{lem:integral}
The integral $x$ coordinates on the following two elliptic curves are
\[
\begin{array}{c|l}
 E_1:\ y^2+y=x^3+12x^2-16 & \{-8,4\}\\
 E_2:\ y^2+y=x^3-12x^2-6x+64 & \{-2,2,16,44,52\}.
\end{array}
\]
\end{lemma}
\begin{proof}[Computer-assisted verification]
The file \texttt{erdos243-elliptic-points.m} constructs precisely these
models in Magma V2.29-10. It explicitly sets both class-group bound maps
to the Minkowski bound, calls \texttt{MordellWeilGroup}, and asserts that
both its rank-proved and full-group-proved booleans are true. It obtains
respectively a group $\Z$, with free generator $(-8,-16)$, and
$\Z/3\oplus\Z$, with free generator $(2,3)$ modulo torsion. It then
calls \texttt{IntegralPoints} with the resulting full free bases and
\texttt{SafetyFactor := 2}, retaining torsion as provided by that routine.
The returned points, up to the involution $y\mapsto-y-1$, are
\[
\begin{array}{c|l}
 E_1&(-8,-16),\ (4,15)\\
 E_2&(-2,4),\ (2,3),\ (16,31),\ (44,248),\ (52,-329).
\end{array}
\]
The script asserts the claimed complete $x$ sets. The actual response,
including all proof flags, is saved in
\texttt{erdos243-elliptic-magma-certified.xml}.
The completeness here uses Magma's implemented integral-point algorithm,
not a bounded-height point search. The relevant documented contracts are
\href{https://magma.maths.usyd.edu.au/magma/handbook/text/1567}{IntegralPoints},
\href{https://magma.maths.usyd.edu.au/magma/releasenotes/2/21/}{the two
Mordell--Weil proof flags}, and
\href{https://magma.maths.usyd.edu.au/magma/handbook/text/416}{the unconditional
Minkowski class-group bounds}. These software computations have not been
formally verified in Lean or reproduced by a second implementation here.
\end{proof}

\begin{theorem}[Centered quadratic profiles]\label{thm:profiles}
No orbit~\eqref{eq:orbit} can have, eventually,
\[
 C_n=A_0 n(n+1)+B_0,\qquad A_0\in\Q_{>0},\quad B_0\in\Q.
\]
The same conclusion holds after any integer shift of the index.
\end{theorem}
\begin{proof}
Clear coefficient denominators by multiplying both states by a fixed
positive integer, so $C_n=A n(n+1)+B$ with $A\in\Z_{>0}$, $B\in\Z$.
Put $G_n=\gcd(C_n,D_n)$. The recurrence gives $G_n\mid G_{n+1}$,
and $G_n$ divides every subsequent $C_j$. Hence $G_n\mid2A$ by
second differences; the chain stabilizes at $g\mid2A$. It also divides
$B$, because every $A n(n+1)$ is a multiple of $2A$. Dividing by $g$
gives primitive states with
\[
 U_n=m\frac{n(n+1)}2+c,\quad m=2A/g\in\Z_{>0},\quad c=B/g\in\Z,
 \qquad \gcd(U_n,V_n)=\gcd(U_n,U_{n+1})=1.
\]
If $c\notin\{-1,1\}$, choose a prime dividing $c$, or any prime
if $c=0$. At every sufficiently late $n$ divisible by $2p$ both
$U_{n-1},U_n$ are divisible by $p$, a contradiction. Thus $c=\pm1$.
If $m$ is odd, $U_n,U_{n+1}$ are both even at $n\equiv1\pmod4$.
So $m$ is even and the primitive profile is
\[
 U_n=k n(n+1)+c,\qquad k\in\Z_{>0},\quad c=\pm1.
\]

For $c=1$, let $f(X)=kX(X+1)+1$, with root $\alpha$ and
$\Delta=k(k-4)$. Equation~\eqref{eq:neighbours} at every root modulo
every good prime forces
$-f(\alpha-1)f(\alpha+1)=-4k$ to be a square in $\Q(\alpha)$ by
Lemma~\ref{lem:cheb}. If $k\ge4$, that field is real, including the
repeated-root case $k=4$, which is impossible for the negative number
$-4k$. For $k=1$ the field is $\Q(\sqrt{-3})$, where $-4$ is not a
square; for $k=2$ it is $\Q(i)$, where $-8$ is not a square.
The remaining $k=3$ has, at $n\equiv1\pmod7$, the four values
$(U_{n-2},U_{n-1},U_n,U_{n+1})\equiv(1,1,0,5)$.
Then $V_n=2$, $V_{n-1}^2=2$, and $V_{n-1}=3$ or $4$ modulo $7$.
The preceding step in~\eqref{eq:quadratic} would require a root of
$x^2+x-V_{n-1}$, with nonsquare discriminant $6$ or $3$ modulo $7$.
This excludes $c=1$.

It remains to treat $c=-1$. Write $f(X)=kX(X+1)-1$ and
$\Delta=k(k+4)$. The discriminant is nonsquare for every positive
integer $k$, since $(k+1)^2<\Delta<(k+2)^2$. The same square test
forces $4k$ to be a square in $K=\Q(\sqrt\Delta)$.
A rational number $h$ is a square in $\Q(\sqrt d)$ only if $h$ or
$h/d$ is a rational square, by expanding $(u+v\sqrt d)^2$.
Consequently
\begin{equation}\label{eq:twofamilies}
 k=s^2\quad(s\ge1),\qquad\hbox{or}\qquad k=s^2-4\quad(s\ge3).
\end{equation}

Here is a second backward step in the square test. At an algebraic root
$\alpha$ of $f$, put $\eta=k(2\alpha+1)$, so
\[
 \eta^2=\Delta,\qquad f(\alpha+j)=kj^2+\eta j.
\]
At a corresponding root modulo a good prime, the state at the root is
$V_0=-f(\alpha+1)$ and its predecessor obeys $V_{-1}^2=4k$.
For each of the two possible predecessor values, the preceding state
must solve a quadratic with discriminant
\begin{equation}\label{eq:backdisc}
 (k-\eta)^2+4(4k-2\eta)V_{-1}.
\end{equation}
These statements are congruences at the root in a finite field; we are
not assuming that the actual integer states are polynomial functions
of $\alpha$.

In the first case of~\eqref{eq:twofamilies}, take $t=\sqrt{s^2+4}$,
so $K=\Q(t)$ and $\eta=st$. The choices $V_{-1}=\sigma 2s$,
$\sigma\in\{-1,1\}$, make~\eqref{eq:backdisc} equal to $s^2b_\sigma$,
where
\[
 b_\sigma=2s^2+4+32\sigma s-(2s+16\sigma)t,
 \qquad
 N_{K/\Q}(b_\sigma)=16(4\sigma s^3+48s^2-63).
\]
In the second case take $t=\sqrt{s^2-4}$, so $\eta=st$ and the
choices are $V_{-1}=\sigma2t$. The discriminants are $t^2b_\sigma$,
where
\[
 b_\sigma=2s^2-4-16\sigma s+(32\sigma-2s)t,
 \qquad
 N_{K/\Q}(b_\sigma)=16(4\sigma s^3-48s^2-24\sigma s+257).
\]
The norm formulas are polynomial identities, also checked symbolically
in the Magma script. None of these $b_\sigma$ is zero, since each norm
divided by $16$ is $1$ modulo $4$.

In either case at least one $b_\sigma$ must be a square in $K$.
Otherwise Lemma~\ref{lem:cheb} supplies a good prime and a root in
$\Fp$ at which both are nonsquares, so neither possible predecessor
can extend one step further backward. This contradicts the integer
orbit, since the selected residue class contains arbitrarily late
indices. The norm of a square is a rational square. Thus in the first
case an integer square equals
$4x^3+48x^2-63$ at $x=\sigma s$; in the second it equals
$4x^3-48x^2-24x+257$. Each square is odd. Writing its square root as
$2y+1$ gives exactly $E_1$ or $E_2$ of Lemma~\ref{lem:integral}.
That complete computation leaves only
\begin{equation}\label{eq:five}
 k\in\{16,64,252,1932,2700\}.
\end{equation}

All five possibilities are excluded by the following small finite-field
certificates. For the indicated prime $p$ and root residue $r$, put
$S_r=\{-f(r+1)\}$, and recursively set
\[
 S_j=\{x\in\Fp^\times:\text{ for some }d\in S_{j+1},\ 
                 x(x+f(j+1))=f(j)d\},\qquad j=r-1,r-2,\ldots .
\]
The successive sets, starting at $S_r$, are
\[
\begin{array}{c|c|c|l}
 k&p&r& S_r,S_{r-1},\ldots\\\hline
 16,2700&11&4&
 \{5\},\{3,8\},\{2,3,4,5\},\{2,5,6,9,10\},\{6,7\},\varnothing\\
 64&13&3&\{8\},\{3,10\},\{10\},\varnothing\\
 252&19&6&\{6\},\{1,18\},\{4,18\},\varnothing\\
 1932&11&5&\{4\},\varnothing
\end{array}
\]
Here all displayed $f(r+1)$ are nonzero. At an arbitrarily late index
in the root class, $V_r=-f(r+1)$ is a unit. Every earlier $V_j$ in
the window is also a unit, since it divides $V_r$. Identity~\eqref{eq:quadratic}
would put their values successively in the displayed sets, which end
empty, a contradiction. The sets can be checked by testing at most
$p-1$ values per row. They are exhaustively checked, both by the
quadratic identity and independently by the original two recurrences,
in \texttt{erdos243-quadratic-certificates.mjs}; the output is retained
in \texttt{erdos243-quadratic-certificates.json}.
This eliminates~\eqref{eq:five} and proves the theorem.
\end{proof}

\begin{theorem}[Quadratic regular-rate irrationality]\label{thm:rate}
If $(a_n)_{n\ge1}$ is a strictly increasing sequence of positive
integers with
\[
 \frac{a_n^2}{a_{n+1}}=1+\frac2n+o(n^{-2}),
\]
then $\sum_{n\ge1}1/a_n$ is irrational.
\end{theorem}
\begin{proof}
Assume the sum is rational, with denominator $q$. Define
$R_n=\sum_{j\ge n}1/a_j$, $D_n=q\prod_{j<n}a_j$, and $C_n=D_nR_n$.
These are positive integer states satisfying~\eqref{eq:orbit}.
The ratio hypothesis implies $a_{n+1}\sim a_n^2$, superpolynomial
growth, and
\[
 R_n=\frac1{a_n}+\frac1{a_{n+1}}+O(a_{n+1}^{-2}),\qquad
 R_n^{-1}=a_n-\frac{a_n^2}{a_{n+1}}+O(a_n^{-1}).
\]
Consequently
\[
 \frac{C_{n+1}}{C_n}=a_n-R_n^{-1}
 =1+\frac2n+o(n^{-2}).
\]
Normalizing by $n(n+1)$, the ratio of consecutive normalized terms is
$1+o(n^{-2})$. Its logarithmic tail is $o(n^{-1})$, so for some
$A>0$,
\[
 C_n=A n(n+1)+o(n),\qquad
 C_{n+1}-C_n=2A(n+1)+o(1).
\]
The consecutive differences are integers, so their own differences
are integers tending to $2A$. Thus $2A=m\in\Z_{>0}$ and
$C_{n+2}-2C_{n+1}+C_n=m$ eventually. The preceding asymptotic then
forces $C_{n+1}-C_n=m(n+1)$ eventually, and hence
\[
 C_n=\frac m2 n(n+1)+B
\]
eventually for a fixed integer $B$. This contradicts
Theorem~\ref{thm:profiles}.
\end{proof}

\paragraph{Unresolved step for the full problem.}
The unrestricted condition $a_{n+1}/a_n^2\to1$ yields relative errors
tending to zero, but supplies neither the regular rate above nor a
polynomial numerator. This proof does not rule out a nonpolynomial
infinite orbit with nonzero errors tending to zero. The full proof or
disproof requested for Erd\H{o}s 243 remains outstanding.
\end{document}
